\section{Conclusiones}
\label{vi:sec:conclusiones}

El resultado principal es la construcción de un operador positivo de masa
sobre la fibra de una partícula persistente, con una especialización
escalar determinada para el electrón central. La positividad se deduce de
la transformación por congruencia del operador basal positivo, con lector
autoadjunto y razón de acción positiva. En la trayectoria electrónica,
los dos lectores producen \((80,54,6)\) y determinan la misma corrección
de retorno. En la coordenatización dimensional adoptada, la evaluación completa da
\[
 m_e^{\beta}=0.5109989506900008086\ldots\ \mathrm{MeV}/c^2.
\]
La cifra reúne el prefactor electrónico y la corrección de retorno; su
contraste se realiza con el valor de referencia y la incertidumbre
indicados en la tabla~\ref{tab:vi-evaluaciones-veinte}. El resultado
escalar específico del centro no sustituye los operadores ni las
condiciones de individualización de las demás colecciones.

La trayectoria que sostiene esa evaluación no es inmóvil: su proyección
celular visita exactamente las veintisiete posiciones definidas por
$c-r\equiv0,\pm1\pmod9$. El registro reversible distingue historias que
comparten un mismo histograma de valores absolutos; la descomposición espectral
de $(+++---)^2$ concentra la potencia en $k=2,6,10$ y reproduce exactamente
$(\Omega,P_6)=(80,6)$. La
condición operatoria sobre la fibra establece además que un transporte TPK
conserva masa si y sólo si conserva la coordenada completa del carácter y del
exponente de retorno, condición equivalente a $[M^\beta,V_F]=0$ en el dominio
estable.

El ambiente regional de \(729\) posiciones tampoco queda como un cardinal
aislado. La biyección por pares de trits
\(\beta:\FF_3^6\to E_9^3\) lo identifica con el cuerpo \(S_0\), mientras
que la completación \(\widehat E^3\) se descompone de forma disjunta en
\(729+243+27+1=1000\) posiciones. La corona completa contiene \(271\)
posiciones, el ápice aporta exactamente \(1/1000\) de la medida uniforme y
la compatibilidad proyectiva conserva la normalización en toda dimensión.
La prueba y las dos figuras del cuerpo exponen por separado la partición
conjuntista y su realización cúbica.

La procedencia de los coeficientes queda determinada por la evolución
ordenada de \(108\) operadores
\(\mathcal U_t=\mathsf{Upd}_t\circ\mathsf{Tra}_t\circ\mathsf{Sel}_t\).
La coordenatización \(\mathcal R_{12}\) extrae del estado terminal el libro
incidencial completo; \(\mathcal E_{108}^{90,120}\) produce la coordenada
de canales; \(\Pi_H\) extrae la firma; y \(H_{12}/4\) obtiene \(K\).
La inversa integral restituye exactamente canales y suma. El registro
periódico, la normalización de \(\alpha\), la escala de acción y los
lectores de masa actúan sobre esas salidas. El mismo origen estructural
se conserva, por tanto, desde las operaciones del núcleo hasta la
evaluación y su comparación metrológica.

\subsection{Identidad persistente, atlas y conjugación}

La persistencia se formula sobre extensiones compatibles; la ruta
conserva su historia observable; la representación determina qué datos de
hoja, orientación, espín, sabor y color distinguen los estados. Esta
identidad precede al operador de masa y permite relacionar sus lecturas
sin reducir la partícula a una coordenada numérica.

El atlas proporciona una clasificación finita que puede recorrerse
íntegramente. Sus \(81\) celdas, \(56\) colecciones, \(13\) multisecciones y
\(324\) condiciones orientadas son cuatro censos de objetos diferentes.
Su organización identifica la excepción central del electrón y las
fibras no centrales. En estas últimas, conservar una multisección o un
operador es una descripción positiva del resultado: precisa qué permanece
equivariante y qué operación adicional produce una lectura espectral.
La individualización física se realiza mediante su representación y su
canal, conservando esa diferencia.

La conjugación tampoco se reduce a una única inversión de signo.
La inversión celular, la reversión de una ruta, la involución sobre
palabras, la conjugación de carga y la holonomía espinorial conservan
dominios y acciones propios. La construcción de Möbius especifica la
fibra y el lazo que generan la inversión de orientación. El retorno a la
misma fase puede dejar memoria distinta, mientras la restitución de una
representación puede exigir una vuelta adicional. Esos resultados
describen cómo persiste la estructura cuando sus observaciones parciales
coinciden.
La conjugación antiunitaria precisa además una distinción dinámica:
la conmutación con el Hamiltoniano invierte el sentido del grupo de
evolución, mientras la preservación de su subespacio real a todo tiempo
exige la relación de anticonmutación correspondiente. La compatibilidad
con el operador de masa conserva su formulación propia.

\subsection{Evaluación másica, composición y ocupación}

La evaluación de masa combina dos niveles. El carácter refinado utiliza
la firma entera y los momentos de los canales; la sección dimensional
compone acción, reloj y coordenatización de unidades. El prefactor electrónico
permanece presente al expresar el electrón respecto a la escala
cronológica. La razón de las secciones de acción actúa en los lectores
de retorno y permite comparar fibras mediante exponentes derivados de
sus registros. La igualdad de los dos lectores centrales ha recibido
una prueba concreta sobre la trayectoria electrónica, con su alcance
diferenciado de los demás perfiles del atlas.
La velocidad del reloj y la velocidad constitutiva coinciden mediante
la identidad entre el producto de los dos canales y el determinante del
operador de respuesta. Al cambiar la coordenatización de unidades se transforman
sus coordenadas y permanece invariante la magnitud dimensional.

La composición de partículas se ha situado antes de la extracción de
la firma. Los registros de enlace conservan las condiciones de frontera,
las hojas y los predecesores sobre los que actúan los funcionales. Por
ello el término de enlace posee un lugar definido en el operador del
compuesto. Mesones y bariones se organizan mediante sus espacios de
constituyentes y sus representaciones cromáticas; las resonancias
requieren además la dinámica de canales abiertos y la definición del
polo. Su anchura y su tiempo de supervivencia son lecturas coordinadas
de esa dinámica.

En el sector antipodal neutral, las dos contracciones escalares del
extractor se anulan exactamente:
\(b_{120}^{\nu}=b_{\zeta}^{\nu}=0\). Esta cancelación no elimina el
registro orientado, pues el subespacio antiinvariante conserva el operador
complejo \(J^2=-I\); tampoco se extiende a las rutas generales ni modifica
las masas o el transporte PMNS. En la prolongación longitudinal, la
comparación entre diez palabras de \(108\) pasos y su concatenación de
\(1080\) pasos incorpora el término de costura explícito. Las nuevas
fronteras no son coronas, de modo que
\(\Delta C_z=0\) y \(\Delta k_{\Delta_4}=\Delta T_z\), con \(-2\) como
valor alcanzable del control local. Así se conservan separadamente la
periodicidad de la representación y la memoria del estado TPK completo.

La completación graduada desarrolla la ocupación sobre los modos
completos ya declarados. La antisimetrización construye las relaciones
CAR y de ellas se deduce la exclusión de Pauli. La simetrización
construye las CCR en el dominio de partículas finitas y hace visible
el término de frontera introducido por un truncamiento bosónico.
La distinción entre igualdad de una masa y coincidencia de un modo
queda expresada por el determinante de Gram de una ocupación doble.
Espín, hoja o color pueden conservar modos independientes aun cuando
un observable parcial tenga el mismo valor.

Las realizaciones orbitales enlazan la acción con las condiciones de
cierre y la geometría de Kepler--Sommerfeld. Sus hipótesis de métrica,
dinámica central y representación se han mantenido junto a las
ecuaciones que las utilizan. En los sectores de fusión y trenzado se
explicitan el anillo de incidencia, sus dimensiones positivas, las
matrices de cambio de asociación y los caracteres de intercambio,
junto con las identidades locales y las condiciones de realización
que utiliza cada uno. Esta distribución permite reconocer qué corresponde a
una construcción algebraica, a una realización dinámica y a una
identificación física.

\Needspace{29\baselineskip}
\subsection{Del espectro al contraste metrológico}

El contraste metrológico completa el recorrido mediante observables
tipados. Las comparaciones explicadas por sectores se acompañan de
los inventarios completos y sus fuentes. Una masa corriente, una masa
de polo, un límite y una anchura conservan sus definiciones, unidades
e incertidumbres. Cada evaluación del modelo se vincula a su registro,
su lector y su coordenatización; cada observación experimental conserva su edición
y su procedencia. Las \(855\) filas documentales de masas y anchuras
constituyen el dominio del catálogo de contraste, distinto del número
de predicciones individualizadas.

La aportación reúne tres determinaciones: el objeto que persiste, la
operación que produce su espectro y la coordenatización en la que ese espectro se
contrasta. La masa resulta de su composición; orientación, memoria y
estructura de fibra siguen disponibles para distinguir estados con una
misma lectura parcial. El caso electrónico realiza esa composición en
una sección escalar explícita, mientras el atlas conserva la diversidad
de los sectores operatorios y de sus realizaciones físicas. Esa relación
entre estructura y magnitud constituye la unidad del resultado.

% BEGIN SINTESIS_ADITIVA_20260922
\par
La composición adicional cierra la dependencia longitudinal y
gravitatoria sin alterar el espectro. La incidencia produce
\(L_\alpha=\alpha_{\rm HMT}^{16}(5759/23040)L_\star\) antes de
construir el reloj; de ahí se obtiene \(R_{108}=L_\alpha\). La energía
\(Q_L=\hbar c_{\rm int}/L_\alpha\) fija conjuntamente
\(m_{\rm clk}\) y el factor cotransformado
\(b_e=E_e^{(0)}/Q_L\), por lo que la escala se cancela en cada masa ya
construida. En las fibras masivas, las dos igualdades conservadas
\(H\Lambda=\hbar c_{\rm int}I\) y
\(\mathcal R_g\Lambda=L_\alpha^2I\) fuerzan de forma única
\(G=c_{\rm int}^3L_\alpha^2/\hbar\). Esta determinación es posterior al
operador de masa; la expresión con \(t_0\) es su corolario circular y no
su generador. El sector nulo conserva su tratamiento no invertible y no
se introduce ninguna corrección decimal dependiente de la especie.
% END SINTESIS_ADITIVA_20260922
